% ===========================================================================
%  sviluppatore/08-qualita/02-diagnosi.tex
% ===========================================================================
\section{Quando qualcosa non va}\label{sez:diagnosi}

Un elenco dei sintomi che capitano davvero, con la causa più probabile
accanto. Ciascuno rimanda al punto della guida che lo spiega.

\begin{center}
\begin{longtable}{@{}L{50mm}L{88mm}@{}}
\toprule
\textbf{Sintomo} & \textbf{Dove guardare} \\
\midrule
\endfirsthead
\toprule
\textbf{Sintomo} & \textbf{Dove guardare} \\
\midrule
\endhead

\rigaTenue
Metà dell'interfaccia è sparita dopo una modifica
  & Un modulo \file{ui\_ext\_*} solleva all'import, e il \cod{try} unico
    impedisce anche agli altri di caricarsi. Nel registro:
    \cod{v2 page extensions not loaded}. \vediSez{ui-ext} \\

Una strategia non compare nell'elenco
  & \file{python tests/diagnostica\_strategie.py}, che stampa l'errore di
    import \emph{visibile} --- nella dashboard finisce in un avviso che
    nessuno legge. \\

\rigaTenue
Un artefatto nuovo non compare
  & Dieci cause possibili, in ordine di frequenza:
    \vediSez{trappole}. \\

Una modalità di pesatura compare ma risponde \cod{400}
  & Ricarica parziale del registro:
    \cod{EnsembleEngine.WEIGHTING\_MODES} non è stato aggiornato.
    \vediSez{trappole} \\

\rigaTenue
Il backtest si ferma alla prima barra con \cod{TypeError}
  & La firma di \cod{generate\_signals} usa nomi di parametro diversi da
    quelli attesi. \vediSez{decisione-esecuzione} \\

Lo Sharpe è straordinario
  & Prima di festeggiare: lookahead. Si legge oltre \cod{t\_decision},
    oppure una soglia è un quantile calcolato sull'intero campione.
    \vediSez{decisione-esecuzione} \\

\rigaTenue
Dal vivo il conto si muove molto più (o molto meno) del backtest
  & La normalizzazione per stream. È il difetto storico del progetto, e la
    suite che lo sorveglia è \file{prova\_normalizzazione\_stream.py}.
    \vediSez{live} \\

Il marketplace è vuoto in locale e pieno sul sito
  & Canale~\num{1}: token mancante o scaduto. Il token è in cache per tutta
    la vita del processo. \vediSez{ponte} \\

\rigaTenue
«Non sei abbonato» per tutti
  & Si sta chiamando il sito con il segreto della dashboard invece che con
    il token dell'utente. \vediSez{ponte} \\

Un download dati fallisce solo su una macchina
  & Canale~\num{2}: la chiave non è arrivata. Le variabili d'ambiente le
    gestisce il sito. \vediSez{due-meta} \\

\rigaTenue
Una pagina mostra «nessun acquisto» a chi ne ha
  & Una lettura fallita trattata come elenco vuoto. \vediSez{pagina-nuova} \\

Un mercato di regime non modula
  & Nessuno dei suoi \cod{driver} è disponibile: lo dice un avviso nel
    registro, non un errore. Serve
    \file{scripts/download\_regime\_drivers.py}. \vediSez{nuovo-regime} \\

\rigaTenue
Le pagine sono tutte vuote all'avvio
  & Non ci sono dati. Non è un errore: non c'è niente da mostrare.
    \vediSez{dati} \\

\percorsoHTTP{/ensemble/core} risponde \cod{404}
  & Gli ensemble preconfigurati sono stati disattivati: esiste solo
    l'ensemble costruito a runtime. \vediSez{data-layer} \\

\bottomrule
\end{longtable}
\end{center}

\subsection{Gli strumenti di diagnosi}

\begin{center}
\begin{tabular}{@{}L{48mm}L{90mm}@{}}
\toprule
\textbf{Strumento} & \textbf{A che serve} \\
\midrule
\rigaTenue
\file{tests/diagnostica\_strategie.py}
  & Per ogni cartella di strategia: importa o dà errore, con il traceback
    intero. Si lancia dalla radice, con lo \emph{stesso} interprete con cui
    si lancia la dashboard --- la differenza fra due ambienti virtuali è una
    causa frequente. \\
\rotta{GET}{/api/marketplace/diagnostica}
  & Verifica la catena della vetrina e dice in quale punto si rompe. \\
\rigaTenue
\rotta{GET}{/api/marketplace/<id>/perche-non-mio}
  & Perché un artefatto non risulta acquistato, distinguendo «non
    comprato» da «non sono riuscito a chiedere». \vediSez{pagamenti} \\
\rotta{GET}{/api/data/cache\_status}
  & Che cosa c'è su disco: numero di file, dimensione, intervallo di date. \\
\rigaTenue
\rotta{GET}{/api/universo}
  & Che cosa il sistema \emph{sa} di avere, contro che cosa ha davvero. La
    differenza fra le due è il registro \file{data/\_universe.json}. \\
\rotta{GET}{/api/streams/orfani}
  & Risultati senza una strategia che li abbia prodotti: restano dopo una
    cancellazione, e sporcano gli elenchi. \\
\bottomrule
\end{tabular}
\end{center}

\begin{pratica}{Il primo posto dove guardare, sempre}
Il registro. Questo sistema è pieno di \cod{try/except} che registrano un
avviso e proseguono, ed è una scelta giusta --- un plugin rotto non deve
impedire l'avvio --- che però sposta l'informazione dal comportamento al
registro. Un sintomo che non ha un errore corrispondente nell'interfaccia ha
quasi sempre una riga nel registro, e quasi sempre quella riga contiene il
nome del modulo colpevole.
\end{pratica}
